
\documentclass{article}
\usepackage{colortbl}
\usepackage{makecell}
\usepackage{multirow}
\usepackage{supertabular}

\begin{document}

\newcounter{utterance}

\centering \large Interaction Transcript for game `dond', experiment `semi\_en', episode 15 with reflexion{-}Qwen3.5{-}9B.
\vspace{24pt}

{ \footnotesize  \setcounter{utterance}{1}
\setlength{\tabcolsep}{0pt}
\begin{supertabular}{c@{$\;$}|p{.15\linewidth}@{}p{.15\linewidth}p{.15\linewidth}p{.15\linewidth}p{.15\linewidth}p{.15\linewidth}}
    \# & $\;$A & \multicolumn{4}{c}{Game Master} & $\;\:$B\\
    \hline

    \theutterance \stepcounter{utterance}  
    & & \multicolumn{4}{p{0.6\linewidth}}{
        \cellcolor[rgb]{0.9,0.9,0.9}{
            \makecell[{{p{\linewidth}}}]{
                \texttt{\tiny{[P1$\langle$GM]}}
                \texttt{You have attempted this same task before and did not succeed. These are the plans you wrote after those attempts.} \\
\texttt{Trial \#1: New plan: I failed because I assumed my proposal of "2 combs, 2 bowls" would be accepted without verifying the other player's actual secret proposal or their valuation of tables. I treated the other player's message "2 tables" as a confirmed agreement rather than a counter{-}proposal that might conflict with mine or be based on a different distribution logic. Next time, I will not immediately submit a secret proposal after the first exchange. Instead, I will send a second public message explicitly confirming the complementary nature of our offers (e.g., "If you take the 2 tables, I will take the 2 combs and 2 bowls, which sums to all available items") to ensure mutual understanding before finalizing the deal with a secret proposal.} \\
\\ 
\texttt{Trial \#2: New plan: I failed because I submitted my secret proposal immediately after the other player stated their desire for the tables, without explicitly confirming that our combined requests (2 tables + 2 combs + 2 bowls) perfectly utilized the entire set of available items. The game requires proposals to be complementary; by not verbally validating that "Your 2 tables + My 2 combs and 2 bowls = All Items" before locking in the secret proposal, I risked a mismatch or a failure to recognize the deal was complete. Next time, after the other player states their item preference, I will send a second public message explicitly stating: "Your request for 2 tables is compatible with my request for 2 combs and 2 bowls, as this covers all available items. I am ready to submit my secret proposal." Only after this confirmation will I submit the formatted secret proposal.} \\
\\ 
\\ 
\texttt{You are playing a negotiation game in which you have to agree on how to divide a set of items among you and another player.} \\
\\ 
\texttt{Rules:} \\
\texttt{(a) You and the other player are given a set of items. Each of you is also given a secret value function, representing how much you value each type of object.} \\
\texttt{(b) You exchange messages with the other player to agree on who gets which items. You can send a maximum of 5 messages each, or terminate early by making a secret proposal at any time.} \\
\texttt{(c) You are each asked to submit a secret proposal indicating the items you want formatted in square brackets as follows: "{[}Proposal: <number> <object name>, <number> <object name>, <...>{]}"} \\
\texttt{(d) If your proposals are complementary, i.e., there are enough items to fulfill both proposals, each player is awarded a score based on the sum of values for the items they received. Otherwise, both of you get zero points.} \\
\texttt{(e) Your goal is to maximize the score you receive.} \\
\\ 
\texttt{Let us start.} \\
\\ 
\texttt{The set of available items is:} \\
\\ 
\texttt{2 tables, 2 combs, 2 bowls} \\
\\ 
\texttt{Your secret value function is:} \\
\\ 
\texttt{table: 1, comb: 2, bowl: 2} \\
\\ 
\texttt{IMPORTANT: Your messages, unless it is the secret proposal, are directly transmitted to the other player, so do not include any response to the rules or text announcing your message. To make a secret proposal, use the indicated format. Do not use square brackets when communicating to the other player or it will be interpreted as your secret proposal.} \\
            }
        }
    }
    & \\ \\

    \theutterance \stepcounter{utterance}  
    & \multicolumn{4}{p{0.6\linewidth}}{
        \cellcolor[rgb]{0.9,0.9,0.9}{
            \makecell[{{p{\linewidth}}}]{
                \texttt{\tiny{[P1$\rangle$GM]}}
                \texttt{I value combs and bowls more than tables (2 vs 1). Since there are 2 of each item, I would like to secure the combs and bowls while letting the other player take the tables.} \\
\\ 
\texttt{I will start by stating my preference for the combs and bowls to gauge their response.} \\
\\ 
\texttt{I want the 2 combs and 2 bowls.} \\
            }
        }
    }
    & & \\ \\

    \theutterance \stepcounter{utterance}  
    & & & \multicolumn{2}{p{0.3\linewidth}}{
        \cellcolor[rgb]{0.9,0.9,0.9}{
            \makecell[{{p{\linewidth}}}]{
                \texttt{\tiny{[GM$|$GM]}}
                \texttt{continue} \\
            }
        }
    }
    & & \\ \\

    \theutterance \stepcounter{utterance}  
    & & & \multicolumn{2}{p{0.3\linewidth}}{
        \cellcolor[rgb]{0.9,0.9,0.9}{
            \makecell[{{p{\linewidth}}}]{
                \texttt{\tiny{[GM$|$GM]}}
                \texttt{I value combs and bowls more than tables (2 vs 1). Since there are 2 of each item, I would like to secure the combs and bowls while letting the other player take the tables.} \\
\\ 
\texttt{I will start by stating my preference for the combs and bowls to gauge their response.} \\
\\ 
\texttt{I want the 2 combs and 2 bowls.} \\
            }
        }
    }
    & & \\ \\

    \theutterance \stepcounter{utterance}  
    & & \multicolumn{4}{p{0.6\linewidth}}{
        \cellcolor[rgb]{0.9,0.9,0.9}{
            \makecell[{{p{\linewidth}}}]{
                \texttt{\tiny{[GM$\rangle$P2]}}
                \texttt{You have attempted this same task before and did not succeed. These are the plans you wrote after those attempts.} \\
\texttt{Trial \#1: New plan: I failed because I misunderstood the critical rule regarding message formatting. I sent "2 tables" as a public message to the other player, which was not a valid negotiation move, and then I used square brackets `{[}{]}` in my secret proposal, which the rules explicitly forbid (stating "Do not use square brackets when communicating to the other player"). The correct strategy is to strictly adhere to the format: send public messages without any brackets or rule commentary, and only use the exact format `{[}Proposal: ...{]}` when making the final secret submission. Next time, I will immediately submit a valid secret proposal in the first turn using the correct format (e.g., `{[}Proposal: 2 comb, 2 bowl{]}`) without sending any preliminary public messages, ensuring no brackets appear in non{-}proposal text.} \\
\\ 
\texttt{Trial \#2: New plan: I failed because I misunderstood the critical rule regarding message formatting. I sent "2 tables" as a public message to the other player, which was not a valid negotiation move, and then I used square brackets `{[}{]}` in my secret proposal, which the rules explicitly forbid (stating "Do not use square brackets when communicating to the other player"). The correct strategy is to strictly adhere to the format: send public messages without any brackets or rule commentary, and only use the exact format `{[}Proposal: ...{]}` when making the final secret submission. Next time, I will immediately submit a valid secret proposal in the first turn using the correct format (e.g., `{[}Proposal: 2 comb, 2 bowl{]}`) without sending any preliminary public messages, ensuring no brackets appear in non{-}proposal text.} \\
\\ 
\\ 
\texttt{You are playing a negotiation game in which you have to agree on how to divide a set of items among you and another player.} \\
\\ 
\texttt{Rules:} \\
\texttt{(a) You and the other player are given a set of items. Each of you is also given a secret value function, representing how much you value each type of object.} \\
\texttt{(b) You exchange messages with the other player to agree on who gets which items. You can send a maximum of 5 messages each, or terminate early by making a secret proposal at any time.} \\
\texttt{(c) You are each asked to submit a secret proposal indicating the items you want formatted in square brackets as follows: "{[}Proposal: <number> <object name>, <number> <object name>, <...>{]}"} \\
\texttt{(d) If your proposals are complementary, i.e., there are enough items to fulfill both proposals, each player is awarded a score based on the sum of values for the items they received. Otherwise, both of you get zero points.} \\
\texttt{(e) Your goal is to maximize the score you receive.} \\
\\ 
\texttt{Let us start.} \\
\\ 
\texttt{The set of available items is:} \\
\\ 
\texttt{2 tables, 2 combs, 2 bowls} \\
\\ 
\texttt{Your secret value function is:} \\
\\ 
\texttt{table: 0, comb: 1, bowl: 4} \\
\\ 
\texttt{IMPORTANT: Your messages, unless it is the secret proposal, are directly transmitted to the other player, so do not include any response to the rules or text announcing your message. To make a secret proposal, use the indicated format. Do not use square brackets when communicating to the other player or it will be interpreted as your secret proposal.} \\
\\ 
\\ 
\texttt{I value combs and bowls more than tables (2 vs 1). Since there are 2 of each item, I would like to secure the combs and bowls while letting the other player take the tables.} \\
\\ 
\texttt{I will start by stating my preference for the combs and bowls to gauge their response.} \\
\\ 
\texttt{I want the 2 combs and 2 bowls.} \\
            }
        }
    }
    & \\ \\

    \theutterance \stepcounter{utterance}  
    & & & \multicolumn{4}{p{0.6\linewidth}}{
        \cellcolor[rgb]{0.9,0.9,0.9}{
            \makecell[{{p{\linewidth}}}]{
                \texttt{\tiny{[GM$\langle$P2]}}
                \texttt{I want the 2 combs and 2 bowls.} \\
            }
        }
    }
     \\ \\

    \theutterance \stepcounter{utterance}  
    & & & \multicolumn{2}{p{0.3\linewidth}}{
        \cellcolor[rgb]{0.9,0.9,0.9}{
            \makecell[{{p{\linewidth}}}]{
                \texttt{\tiny{[GM$|$GM]}}
                \texttt{continue} \\
            }
        }
    }
    & & \\ \\

    \theutterance \stepcounter{utterance}  
    & & & \multicolumn{2}{p{0.3\linewidth}}{
        \cellcolor[rgb]{0.9,0.9,0.9}{
            \makecell[{{p{\linewidth}}}]{
                \texttt{\tiny{[GM$|$GM]}}
                \texttt{I want the 2 combs and 2 bowls.} \\
            }
        }
    }
    & & \\ \\

    \theutterance \stepcounter{utterance}  
    & & \multicolumn{4}{p{0.6\linewidth}}{
        \cellcolor[rgb]{0.9,0.9,0.9}{
            \makecell[{{p{\linewidth}}}]{
                \texttt{\tiny{[P1$\langle$GM]}}
                \texttt{I want the 2 combs and 2 bowls.} \\
            }
        }
    }
    & \\ \\

    \theutterance \stepcounter{utterance}  
    & \multicolumn{4}{p{0.6\linewidth}}{
        \cellcolor[rgb]{0.9,0.9,0.9}{
            \makecell[{{p{\linewidth}}}]{
                \texttt{\tiny{[P1$\rangle$GM]}}
                \texttt{I see that you also want the 2 combs and 2 bowls. Since there are only 2 combs and 2 bowls available, we cannot both have them. We need to split these items or agree on a different distribution.} \\
\\ 
\texttt{Given my valuation (comb: 2, bowl: 2, table: 1) and yours (implied similar or unknown), I need to propose a split that works for both. If you take the tables, I can take the combs and bowls. But you said you want the combs and bowls.} \\
\\ 
\texttt{Let's try a compromise: I will take 1 comb and 1 bowl, and you can take the other comb, the other bowl, and the 2 tables. This way, we both get some of the high{-}value items, and the tables are covered.} \\
\\ 
\texttt{Wait, let me re{-}evaluate. If I take 1 comb and 1 bowl, and you take 1 comb, 1 bowl, and 2 tables, that uses all items. My score would be 2+2=4. Your score would depend on your valuation, but if you value tables less than combs/bowls, this might work.} \\
\\ 
\texttt{Alternatively, if you insist on all combs and bowls, I must take the tables. But I value tables less. Let's see if we can agree on me taking the tables and you taking the combs and bowls. But you already said you want the combs and bowls, implying you don't want the tables.} \\
\\ 
\texttt{Actually, the most logical split is:} \\
\texttt{{-} You: 2 tables} \\
\texttt{{-} Me: 2 combs, 2 bowls} \\
\\ 
\texttt{This uses all items. If you value tables highly, you might prefer this. If you value combs/bowls more, you might prefer the other way. But since you stated you want the combs and bowls, I should propose that I take the tables, and you take the combs and bowls.} \\
\\ 
\texttt{But wait, the rules say we need complementary proposals. If I propose {[}Proposal: 2 tables{]}, and you propose {[}Proposal: 2 combs, 2 bowls{]}, that's complementary.} \\
\\ 
\texttt{However, I need to confirm this with you first. Let me send a public message to confirm.} \\
\\ 
\texttt{If you take the 2 tables, I will take the 2 combs and 2 bowls, which sums to all available items. This seems like a fair split given the valuations.} \\
\\ 
\texttt{I am ready to submit my secret proposal if you agree.} \\
\\ 
\texttt{But I cannot submit the secret proposal yet; I need to confirm with you first.} \\
\\ 
\texttt{Let me try to propose a split where I take the tables, and you take the combs and bowls.} \\
\\ 
\texttt{My proposal: I will take the 2 tables.} \\
\\ 
\texttt{But I need to make sure you agree. Let me send a message.} \\
\\ 
\texttt{I propose that you take the 2 tables,} \\
            }
        }
    }
    & & \\ \\

    \theutterance \stepcounter{utterance}  
    & & & \multicolumn{2}{p{0.3\linewidth}}{
        \cellcolor[rgb]{0.9,0.9,0.9}{
            \makecell[{{p{\linewidth}}}]{
                \texttt{\tiny{[GM$|$GM]}}
                \texttt{proposal} \\
            }
        }
    }
    & & \\ \\

    \theutterance \stepcounter{utterance}  
    & & & \multicolumn{2}{p{0.3\linewidth}}{
        \cellcolor[rgb]{0.9,0.9,0.9}{
            \makecell[{{p{\linewidth}}}]{
                \texttt{\tiny{[GM$|$GM]}}
                [2, 0, 0]
            }
        }
    }
    & & \\ \\

    \theutterance \stepcounter{utterance}  
    & & \multicolumn{4}{p{0.6\linewidth}}{
        \cellcolor[rgb]{0.9,0.9,0.9}{
            \makecell[{{p{\linewidth}}}]{
                \texttt{\tiny{[GM$\rangle$P2]}}
                \texttt{The other player has submitted a secret proposal. It is now your turn to submit a proposal of your own.} \\
            }
        }
    }
    & \\ \\

    \theutterance \stepcounter{utterance}  
    & & & \multicolumn{4}{p{0.6\linewidth}}{
        \cellcolor[rgb]{0.9,0.9,0.9}{
            \makecell[{{p{\linewidth}}}]{
                \texttt{\tiny{[GM$\langle$P2]}}
                \texttt{{[}Proposal: 2 comb, 2 bowl{]}} \\
            }
        }
    }
     \\ \\

    \theutterance \stepcounter{utterance}  
    & & & \multicolumn{2}{p{0.3\linewidth}}{
        \cellcolor[rgb]{0.9,0.9,0.9}{
            \makecell[{{p{\linewidth}}}]{
                \texttt{\tiny{[GM$|$GM]}}
                \texttt{proposal} \\
            }
        }
    }
    & & \\ \\

    \theutterance \stepcounter{utterance}  
    & & & \multicolumn{2}{p{0.3\linewidth}}{
        \cellcolor[rgb]{0.9,0.9,0.9}{
            \makecell[{{p{\linewidth}}}]{
                \texttt{\tiny{[GM$|$GM]}}
                [0, 2, 2]
            }
        }
    }
    & & \\ \\

    \theutterance \stepcounter{utterance}  
    & & & \multicolumn{2}{p{0.3\linewidth}}{
        \cellcolor[rgb]{0.9,0.9,0.9}{
            \makecell[{{p{\linewidth}}}]{
                \texttt{\tiny{[GM$|$GM]}}
                [[2, 0, 0], [0, 2, 2]]
            }
        }
    }
    & & \\ \\

\end{supertabular}
}

\end{document}
